\documentclass[titlepage,a4paper]{jsarticle}
\usepackage{../../../sty/import}% 各種パッケージインポート
\usepackage{../../../sty/title}% タイトルページの変更
\usepackage{../../../sty/answergrid}
%% タイトルページの変数
% レポートタイトル
\title{心理学特論第9回目出席票}
% 提出日
\expdate{\today}
% 科目名
\subject{心理学特論}
% 分野
\class{情報経営システム工学分野}
% 学年
\grade{M1}
% 学籍番号
\mynumber{24336488}
% 記述者
\author{本間三暉}

\begin{document}
% titleページ作成
\maketitle
\section*{keyword\#1}
防衛機制
\section*{keyword\#2}
心理的安全性

\section*{心理学特論第10回テーマの「アドリブで行こう！心を守るコミュニケーション」で取り上げた「雑談用の話のネタ＝雑談のネタ帳」を作成しましょう。
実際に、これからの話のネタとして使えるようにリアルに、具体的に書くのがポイントです。}
\section{マイブーム（スライド31対応）→これまでに何だかやけにはまったことなどがあれば簡潔に書いて下さい。}
私のマイブームは料理である．大学生になり一人暮らしを始めてから，自炊をする機会が増えたことがきっかけで料理に興味を持つようになった．
最初は食費を節約するためという理由もあったが，今では趣味としての側面の方が大きいと感じている．

私は食べることが好きなので，普段からYouTubeやInstagramで料理動画を見ることが多い．特にSNSでは次々と美味しそうな料理が流れてくるため，「これは作れそうだな」と思ったものを保存しておき，休日などに実際に作ってみることがある．
また，「今日はパスタが食べたい」のように作りたい料理が決まっているときは，クックパッドでレシピを探して参考にしている．

これまでにカレーやパスタ，チャーハンなどの定番料理はもちろん，チャーシューや煮卵など時間のかかる料理にも挑戦したことがある．
長時間煮込んだり味を染み込ませたりする工程は大変だが，完成したときの達成感が大きく，料理の面白さを感じる瞬間でもある．

最近作った料理の中で特に美味しかったのは，クリームチーズと牛乳，ケチャップを使ったトマトクリームパスタである．
最初は「本当にこれで美味しくなるのか」と半信半疑だったが，実際に作ってみると想像以上に濃厚で，店で食べるような味に仕上がった．
それ以来，何度か作るお気に入りのレシピになっている．

一方で失敗した経験も多い．例えばパスタはワンパンパスタという，フライパン一つで作る方法をよく使うのだが，始めたばかりの頃は水の量の調整が難しかった．
「少し水が少ないから足そう」と思ったら今度は多すぎたりして，なかなかうまく作れなかった記憶がある．
しかし，何度か作るうちに感覚が分かるようになり，自分でも成長を感じられた．

以前には「暗殺者のパスタ」という少し変わった名前のパスタを作り，友人に食べてもらったこともある．
少し緊張したが，「美味しい」と言ってもらうことができ，自分が作った料理を誰かに喜んでもらえる嬉しさを感じた．

最近はSNSで見るような，大きな肉をじっくり焼く本格的なBBQにも興味を持っている．
日本ではあまり見かけないスタイルだが，いつか友人たちと挑戦してみたいと考えている．
料理は食べる楽しさだけでなく，作る楽しさや新しいことに挑戦する楽しさもあるため，これからも様々な料理に挑戦していきたい．

\section{ご当地自慢など→長岡以外でも結構です。どこの場所かは明記して下さい。}
私が自慢したいご当地は長岡である．長岡は正直なところ，車がないと遊びに行ける場所はあまり多くない．
しかし，人が口に入れるものに関しては本当にレベルが高いと思う．
特におすすめしたいのはラーメンと日本酒である．

私は中学校まで新潟市で過ごし，高専進学を機に長岡へ来た．
長岡に来る前は「ラーメンが有名らしい」という話は聞いていたものの，「美味しいかもしれないけど，わざわざ食べに行くほどではないだろう」と思っていた．
しかし実際に住み始めると，周囲の友人たちに連れられて様々な店を巡るようになり，気が付けば週に1回以上ラーメンを食べる生活になっていた．
最近では食べ過ぎのせいか，少しお腹周りが気になっているほどである．

長岡はラーメン激戦区として知られており，長岡生姜醤油ラーメンのようなあっさりしたものから，濃厚なラーメン，二郎系インスパイアのようながっつりしたラーメンまで様々な種類を楽しむことができる．
私自身は生姜醤油ラーメンも好きだが，特別それだけが好きというわけではなく，色々な種類のラーメンを食べ歩くのが好きである．
中でも「麺屋 松」や「しょうがの海」は自分だけでなく周囲の友人からの評価も高く，長岡に来た際にはぜひ食べてみてほしい店である．

また，長岡は日本酒も有名である．
私はお酒を飲むことが好きで，日本酒もよく飲むが，長岡に来てから地酒を飲む機会が増えた．
特に吉乃川は有名で，長岡を代表する日本酒の一つだと思う．
長岡技術科学大学には，日本酒好きの学生が集まる「しゅがく」というサークルも存在するほどである．

さらに学内では「酒陣」というイベントが開催されることがあり，複数の酒蔵の方が大学に来て様々な日本酒を試飲できる．
私も参加したことがあるが，非常に楽しく，気に入った日本酒を購入した結果，気が付けば酒瓶を6本ほど買ってしまった．
それくらい魅力的な日本酒が多いということだと思う．

長岡は観光地としては地味かもしれないが，ラーメンや日本酒など美味しいものに恵まれた街である．
もし長岡を訪れる機会があれば，ぜひ美味しいラーメンと日本酒を楽しんでほしいと思う．

\section{いたわりワードを考える→自分が他の人から言われたら心地よく感じるであろう、相手をいたわる言葉（セリフ）を、どのようなシチュエーションかも簡潔に添えて考えてみましょう。}
シチュエーション：
就職活動や研究で思うような結果が出ず，自分を責めているとき．

いたわりワード：
「結果はどうあれ，今までやってきたことは無駄じゃないと思うよ．」

私はうまくいかなかったときに，「もっとできたのではないか」と考えてしまうことがある．
そのため，結果だけではなく，そこまで取り組んできた過程を認めてもらえる言葉をかけられると嬉しい．
また，自分の努力を理解しようとしてくれていることも伝わるため，安心して話すことができると思う．
\newpage
\section{心理学特論課題サイコロワーク7回答記入欄}
\AnswerGridFilled{
  {生活ストレス}
  {楽観主義}
  {予測可能}
  {公的機関}
  {PFA}
  {問題解決}
  {達成可能}
  {成功体験}
  {コントロール}
  {親や保護者}
  {被災者に近づき，活動を始める}
  {安全と安心}
  {安定化}
  {情報を集める}
  {現実的な問題の解決を助ける}
  {周囲の人々との関わりを促進する}
  {対処に役立つ情報}
  {紹介と引き継ぎ}
  {早期支援を提供する}
  {時間を調節}
}
\end{document}